% Part: first-order-logic
% Chapter: introduction
% Section: satisfaction

\documentclass[../../../include/open-logic-section]{subfiles}

\begin{document}

\olfileid[pt]{fol}{int}{sat}

\olsection{Satisfação}

Já podemos avançar para a semântica da lógica de primeira ordem assim
que sabemos o que são !!{formula}s: aqui, a definição básica é a de
!!a{structure}. Para nossa linguagem simples, uma !!{structure}
~$\Struct M$ tem apenas três componentes: um conjunto não vazio
$\Domain{M}$ chamado de \emph{!!{domain}}, o que $\Obj a$ designa
em~$\Struct M$, e do que $\Obj P$ é verdadeiro em~$\Struct M$.  O
objeto designado por~$\Obj a$ é denotado por~$\Assign{\Obj a}{M}$ e o
conjunto de coisas das quais $\Obj P$ é verdadeiro
por~$\Assign{\Obj P}{M}$. !!^a{structure}~$\Struct{M}$ consiste
apenas nestas três coisas: $\Domain{M}$, $\Assign{\Obj a}{M} \in
\Domain{M}$ e $\Assign{\Obj P}{M} \subseteq \Domain{M}$. O caso geral
será mais complicado, pois haverá muitos !!{predicate}s e
!!{constant}s, os !!{constant}s podem ter mais de um lugar, e haverá
também !!{function}s.

Isso basta para dar uma definição de satisfação para !!{formula}s que
não contêm !!{variable}s.  A ideia é dar uma definição indutiva que
espelha a maneira como definimos !!{formula}s. Especificamos quando
uma fórmula atômica é satisfeita em~$\Struct{M}$, e então quando, por
exemplo, $\lnot !A$ é satisfeita em~$\Struct{M}$ com base em se $!A$
é ou não satisfeita em~$\Struct{M}$. Por exemplo, poderíamos definir:
\begin{enumerate}
  \item $\Atom{\Obj P}{\Obj a}$ é satisfeita em~$\Struct{M}$ se, e
  somente se, $\Assign{\Obj a}{M} \in \Assign{\Obj P}{M}$.
  \item $\lnot !A$ é satisfeita em~$\Struct{M}$ se, e somente se, $!A$
  não é satisfeita em~$\Struct{M}$.
  \item $(!A \land !B)$ é satisfeita em~$\Struct{M}$ se, e somente se,
  $!A$ é satisfeita em~$\Struct{M}$ e $!B$ também é satisfeita
  em~$\Struct{M}$.
\end{enumerate}
Digamos que $\Domain{M} = \{0, 1, 2\}$, $\Assign{\Obj a}{M} = 1$ e
$\Assign{\Obj P}{M} = \{1, 2\}$. Esta definição nos diria que
$\Atom{\Obj P}{\Obj a}$ é satisfeita em~$\Struct{M}$ (pois $\Assign{\Obj
a}{M} = 1 \in \{1,2\} = \Assign{\Obj P}{M}$). Nos diria ainda que
$\lnot \Atom{\Obj P}{\Obj a}$ não é satisfeita em~$\Struct{M}$, e que,
por sua vez, $\lnot\lnot \Atom{\Obj P}{\Obj a}$ é, e que $(\lnot
\Atom{\Obj P}{\Obj a} \land \Atom{\Obj P}{\Obj a})$ não é satisfeita, e
assim por diante.

A dificuldade surge quando queremos dar uma definição para os
quantificadores: gostaríamos de dizer algo como: ``$\lexists[\Obj
v_0][\Atom{\Obj P}{\Obj v_0}]$ é satisfeita se, e somente se,
$\Atom{\Obj P}{\Obj v_0}$ é satisfeita.'' Mas a !!{structure}
~$\Struct{M}$ não nos diz o que fazer com !!{variable}s. O que
realmente queremos dizer é que $\Atom{\Obj P}{\Obj v_0}$ é satisfeita
\emph{para algum valor de~$\Obj v_0$}.  Para tornar isso preciso,
precisamos de uma maneira de atribuir !!{element}s de~$\Domain{M}$ não
apenas a $\Obj a$, mas também a~$\Obj v_0$. Para isso, introduzimos
\emph{atribuições de variáveis}. Uma atribuição de variáveis é
simplesmente uma função~$s$ que mapeia !!{variable}s para
!!{element}s de~$\Domain{M}$ (em nosso exemplo, para um entre $1$,
$2$ ou~$3$).  Como não sabemos de antemão quais !!{variable}s podem
aparecer em uma !!{formula}, não podemos limitar a quais !!{variable}s
$s$ atribui valores. A solução simples é exigir que $s$ atribua valores
a \emph{todas} as !!{variable}s~$\Obj v_0$, $\Obj v_1$, \dots\@
Usaremos apenas aquelas de que precisarmos.

Em vez de definir a satisfação de !!{formula}s apenas em relação a uma
!!{structure}, definiremos em relação a uma !!{structure}~$\Struct{M}$
\emph{e} uma atribuição de variáveis~$s$, e escreveremos
$\Sat{M}{!A}[s]$ por brevidade. Nossa definição agora incluirá uma
cláusula adicional para lidar com !!{formula}s atômicas contendo
!!{variable}s:
\begin{enumerate}
  \item $\Sat{M}{\Atom{\Obj P}{\Obj a}}[s]$ se, e somente se,
  $\Assign{\Obj a}{M} \in \Assign{\Obj P}{M}$.
  \item $\Sat{M}{\Atom{\Obj P}{\Obj v_i}}[s]$ se, e somente se,
  $s(\Obj v_i) \in \Assign{\Obj P}{M}$.
  \item $\Sat{M}{\lnot !A}[s]$ se, e somente se, não $\Sat{M}{!A}[s]$.
  \item $\Sat{M}{(!A \land !B)}[s]$ se, e somente se, $\Sat{M}{!A}[s]$ e
  $\Sat{M}{!B}[s]$.
\end{enumerate}
Ok, isso resolve um problema: agora podemos dizer quando $\Struct{M}$
satisfaz $\Atom{\Obj P}{\Obj v_0}$ para o valor~$s(\Obj v_0)$. Para
obter a definição correta para $\lexists[\Obj v_0][\Atom{\Obj P}{\Obj
v_0}]$ temos que fazer mais uma coisa: queremos que
$\Sat{M}{\lexists[\Obj v_0][\Atom{\Obj P}{\Obj v_0}]}[s]$ se, e somente
se, $\Sat{M}{\Atom{\Obj P}{\Obj v_0}}[s']$ para \emph{alguma} maneira
$s'$ de atribuir um valor a~$\Obj v_0$. Mas o valor atribuído a~$\Obj
v_0$ não precisa necessariamente ser o valor que $s(\Obj v_0)$ escolhe.
Introduziremos uma notação para isso: se $m \in \Domain{M}$, então
deixamos $\Subst{s}{m}{\Obj v_0}$ ser a atribuição que é igual a~$s$
(para todas as !!{variable}s diferentes de~$\Obj v_0$), exceto que a
$\Obj v_0$ atribui~$m$. Agora nossa definição pode ser:
\begin{enumerate}\setcounter{enumi}{4}
  \item $\Sat{M}{\lexists[\Obj v_i][!A]}[s]$ se, e somente se,
  $\Sat{M}{!A}[\Subst{s}{m}{\Obj v_i}]$ para algum~$m \in \Domain{M}$.
\end{enumerate}
Funciona? Digamos que $s(\Obj v_i) = 0$ para todo~$i \in \Nat$.
$\Sat{M}{\lexists[\Obj v_0][\Atom{\Obj P}{\Obj v_0}]}[s]$ se, e somente
se, há um $m \in \Domain{M}$ tal que $\Sat{M}{\Atom{\Obj P}{\Obj
v_0}}[\Subst{s}{m}{\Obj v_0}]$. E há: podemos escolher $m = 1$ ou
$m = 2$. Note que isso é verdadeiro mesmo que o valor~$s(\Obj v_0)$
atribuído a~$\Obj v_0$ por $s$---neste caso, $0$---não faça o
trabalho. Temos $\Sat{M}{\Atom{\Obj P}{\Obj v_0}}[\Subst{s}{1}{\Obj
v_0}]$, mas não $\Sat{M}{\Atom{\Obj P}{\Obj v_0}}[s]$.

Se isso parece confuso e trabalhoso: é. Mas a complexidade adicional é
necessária para dar uma definição precisa e indutiva de satisfação para
todas as !!{formula}s, e precisamos de algo assim para definir
precisamente as noções semânticas.  Há outras maneiras de fazer isso,
mas todas são igualmente (in)elegantes.

\end{document}
